\subsection{广群的逆向系统}

定义1. 关于指标集I的广群逆向系统是一个集合的逆向系统 \((\mathbf{E}_{\alpha}, f_{\alpha\beta})\) ，其中每个 \(E_{\alpha}\) 具有一个广群结构，且每个 \(f_{\alpha\beta}\) 都是广群同态。

设 \((\mathbf{E}_{\alpha}, f_{\alpha \beta})\) 是一个其律以乘法记法书写的广群逆向系统。逆极限集 \(\mathbf{E} = \varprojlim \mathbf{E}_{\alpha}\) 是积广群 \(\prod_{\alpha \in I} \mathbf{E}_{\alpha}\) 的子集，由这样的族 \((x_{\alpha})_{\alpha \in I}\) 组成：满足 \(x_{\alpha} = f_{\alpha \beta}(x_{\beta})\) （对 \(\alpha \leqslant \beta\) ）。若 \((x_{\alpha})\) 与 \((y_{\alpha})\) 属于 E，则对 \(\alpha \leqslant \beta\) 有 \(x_{\alpha} = f_{\alpha \beta}(x_{\beta})\) 与 \(y_{\alpha} = f_{\alpha \beta}(y_{\beta})\) ，因此 \(x_{\alpha}y_{\alpha} = f_{\alpha \beta}(x_{\beta})f_{\alpha \beta}(y_{\beta}) = f_{\alpha \beta}(x_{\beta}y_{\beta})\) ；因此 E 是 \(\prod_{\alpha \in I} \mathbf{E}_{\alpha}\) 的一个子广群。将给 E 赋予由 \(\prod_{\alpha \in I} \mathbf{E}_{\alpha}\) 上的律诱导的律；如此得到的广群称为诸广群 \(E_{\alpha}\) 的逆极限广群。它满足以下泛性质：

\begin{enumerate}
\def\labelenumi{(\alph{enumi})}
\item
  对于所有 \(\alpha \in \mathbf{I}\) ，典范映射 \(f_{\alpha}\) （它把 E 映入 \(\mathrm{E}_{\alpha}\) ）是广群同态，且 \(\mathrm{E}_{\alpha} \cdot f_{\alpha} = f_{\alpha \beta} \circ f_{\beta}\) 对 \(\alpha \leqslant \beta\) 成立。
\item
  假设给定一个广群 F 以及同态 \(u_{\alpha} \colon \mathrm{F} \to \mathrm{E}_{\alpha}\) ，满足 \(u_{\alpha} = f_{\alpha \beta} \circ u_{\beta}\) （对 \(\alpha \leqslant \beta\) ）。存在一个且仅一个同态 \(u \colon \mathrm{F} \to \mathrm{E}\) ，使得 \(u_{\alpha} = f_{\alpha} \circ u\) 对所有 \(\alpha \in \mathrm{I}\) 成立（即 \(x \mapsto u(x) = (u_{\alpha}(x))_{\alpha \in \mathrm{I}}\) ）。
\end{enumerate}

如果诸广群 \(E_{\alpha}\) 是结合的（相应地，交换的），那么 E 也是。假设每个广群 \(E_{\alpha}\) 有一个单位元 \(e_{\alpha}\) ，且诸同态 \(f_{\alpha\beta}\) 是保单位的。则 \(e = (e_{\alpha})_{\alpha \in I}\) 属于 E，因为有 \(e_{\alpha} = f_{\alpha\beta}(e_{\beta})\) （对 \(\alpha \leqslant \beta\) ），且它是广群 E 的单位元；用上述记号，诸同态 \(f_{\alpha}\) 是保单位的，且若诸 \(u_{\alpha}\) 保单位，则 u 也保单位。进一步，E 的元素 \(x = (x_{\alpha})_{\alpha \in I}\) 可逆当且仅当每个 \(x_{\alpha}\) 在相应的广群 \(E_{\alpha}\) 中可逆，且 \(x^{-1} = (x_{\alpha}^{-1})_{\alpha \in I}\) ；这由公式 \(f_{\alpha\beta}(x_{\beta}^{-1}) = f_{\alpha\beta}(x_{\beta})^{-1} = x_{\alpha}^{-1}\) （对 \(\alpha \leqslant \beta\) ）得出。

由这些评述可以推断，若诸广群 \(E_{\alpha}\) 是幺半群（相应地，群），且诸 \(f_{\alpha\beta}\) 是幺半群同态，则广群 E 是幺半群（相应地，群）。这时我们说这是一个幺半群（相应地，群）的逆向系统。泛性质立即推广到这种情形。

留给读者定义环的逆向系统 \((\mathrm{E}_{\alpha}, f_{\alpha \beta})\) 并验证 \(\mathbf{E} = \varprojlim \mathbf{E}_{\alpha}\) 是积环 \(\prod_{\alpha \in I} \mathbf{E}_{\alpha}\) 的子环，称为诸环 \(\mathbf{E}_{\alpha}\) 的逆极限环；可以验证，泛性质在此情形下仍然成立。

设 \(\mathfrak{E} = (\mathrm{E}_{\alpha}, f_{\alpha\beta})\) 与 \(\mathfrak{E}' = (\mathrm{E}_{\alpha}', f_{\alpha\beta}')\) 是广群（相应地，幺半群、群、环）的两个逆向系统。从 \(\mathfrak{E}\) 到 \(\mathfrak{E}'\) 的一个同态是一个映射的逆向系统 \((u_{\alpha})_{\alpha \in I}\) （ \(u_{\alpha}: \mathrm{E}_{\alpha} \to \mathrm{E}_{\alpha}'\) ），使得每个 \(u_{\alpha}\) 都是同态。在这些条件下，映射 \(u = \varprojlim u_{\alpha}\) 从 \(\varprojlim \mathrm{E}_{\alpha}\) 到 \(\varprojlim \mathrm{E}_{\alpha}'\) 是一个同态（参见《集合论》III，§ 7，第 2 号）。

